\clearpage\pagestyle{otMain}\setlength{\otBlocEspacement}{2.5mm}\fontsize{11}{15}\selectfont\corrigesChapitre{13}{Aires et angles : comparer}
\begin{blocLivre}{cor-13-1}\begin{corrige}{1}{Qui a la plus grande aire ?}\demarche{Compter les carreaux de chaque figure.}
La première couvre 8 carreaux, la seconde 7 : la première a $\res{\text{la plus grande aire}}$.
\label{fig-cor-13-1}\figureGeo[corrige]{aa-1}
\end{corrige}\end{blocLivre}
\begin{blocLivre}{cor-13-2}\begin{corrige}{2}{Mesurer en carreaux}\demarche{Découper la figure en morceaux simples.}
La figure se découpe en un rectangle de 4 carreaux et un carré de 3 carreaux au-dessous : $4+3=\res7$ carreaux.
\label{fig-cor-13-2}\figureGeo[corrige]{aa-2}
\end{corrige}\end{blocLivre}
\begin{blocLivre}{cor-13-3}\begin{corrige}{3}{Formes différentes, aires égales}\demarche{Compter puis comparer.}
Chacune couvre $\res6$ carreaux : même aire, formes différentes. Déplacer un carreau ne change pas l’aire.
\label{fig-cor-13-3}\figureGeo[corrige]{aa-3}
\end{corrige}\end{blocLivre}
\begin{blocLivre}{cor-13-4}\begin{corrige}{4}{En centimètres carrés}\demarche{Compter les carreaux rangée par rangée.}
Le rectangle couvre 4 rangées de 3 carreaux : $4\times3=\res{12\ \mathrm{cm}^2}$.
\label{fig-cor-13-4}\figureGeo[corrige]{aa-4}
\end{corrige}\end{blocLivre}
\begin{blocLivre}{cor-13-5}\begin{corrige}{5}{Nommer un angle}\demarche{Repérer d’abord le sommet.}
C’est l’angle $\widehat{ABC}$ : son sommet est $\res B$, ses côtés passent par A et par C. La lettre du sommet s’écrit au milieu, sous le chapeau.
\label{fig-cor-13-5}\figureGeo[corrige]{aa-5}
\end{corrige}\end{blocLivre}
\begin{blocLivre}{cor-13-6}\begin{corrige}{6}{Ranger des angles}\demarche{Comparer les écartements, pas les longueurs.}
En superposant un calque, l’ordre est $\res{\widehat{B}<\widehat{A}<\widehat{C}}$. Les longueurs des côtés dessinés ne comptent pas.
\label{fig-cor-13-6}\figureGeo[corrige]{aa-6}
\end{corrige}\end{blocLivre}
\begin{blocLivre}{cor-13-7}\begin{corrige}{7}{Chercher les angles droits}\demarche{Poser le coin de l’équerre sur chaque sommet.}
L’équerre s’ajuste exactement en A et en D : la figure possède $\res2$ angles droits, marqués d’un petit carré.
\label{fig-cor-13-7}\figureGeo[corrige]{aa-7}
\end{corrige}\end{blocLivre}
\begin{blocLivre}{cor-13-8}\begin{corrige}{8}{Des moitiés de carreaux}\demarche{Associer les demi-carreaux deux par deux.}
On compte 4 carreaux entiers et 2 demi-carreaux : $4+1=\res5$ carreaux. Deux moitiés se réunissent en un carreau entier.
\label{fig-cor-13-8}\figureGeo[corrige]{aa-8}
\end{corrige}\end{blocLivre}
\begin{blocLivre}{cor-13-9}\begin{corrige}{9}{Même tour, autre surface}\demarche{Calculer séparément contour et surface.}
Périmètres : $2\times(5+1)=12$ et $2\times(3+3)=12$ : égaux. Aires : $5\times1=5$ carreaux contre $3\times3=9$ carreaux : le carré a $\res{\text{la plus grande aire}}$. Même périmètre ne signifie pas même aire.
\label{fig-cor-13-9}\figureGeo[corrige]{aa-9}
\end{corrige}\end{blocLivre}
\begin{blocLivre}{cor-13-10}\begin{corrige}{10}{L’angle du calque}\demarche{Superposer sommet et un côté avant de conclure.}
Le calque coïncide, sommet sur sommet et côtés sur côtés, avec l’angle $\res{\widehat{E}}$ : deux angles superposables sont égaux, même si leurs côtés dessinés ont des longueurs différentes.
\label{fig-cor-13-10}\figureGeo[corrige]{aa-10}
\end{corrige}\end{blocLivre}
\begin{blocLivre}{cor-13-11}\begin{corrige}{11}{Les deux parcelles}\demarche{Compter carreaux entiers puis moitiés.}
P couvre 12 carreaux et Q en couvre 10, moitiés regroupées. P a la plus grande aire : $\res{12\ \mathrm m^2}$ contre 10 $\mathrm m^2$.
\label{fig-cor-13-11}\figureGeo[corrige]{aa-11}
\end{corrige}\end{blocLivre}
\begin{blocLivre}{cor-13-12}\begin{corrige}{12}{La frise de triangles}\demarche{Convertir les triangles en carreaux.}
On compte 14 triangles ; deux triangles font un carreau : $14:2=\res7$ carreaux.
\label{fig-cor-13-12}\figureGeo[corrige]{aa-12}
\end{corrige}\end{blocLivre}
\begin{blocLivre}{cor-13-13}\begin{corrige}{13}{Le côté qui double}\demarche{Compter les carreaux des deux carrés.}
Petit carré : $2\times2=\res4$ carreaux. Grand carré : $4\times4=\res{16}$ carreaux. L’aire est multipliée par $16:4=4$ : elle $\res{\text{a quadruplé}}$, et non doublé.
\label{fig-cor-13-13}\figureGeo[corrige]{aa-13}
\end{corrige}\end{blocLivre}
\begin{blocLivre}{cor-13-14}\begin{corrige}{14}{Les angles du triangle}\demarche{Tester chaque sommet avec l’équerre.}
L’angle en A est $\res{\text{droit}}$ : l’équerre s’y ajuste. Les angles en B et en C sont $\res{\text{plus petits}}$ qu’un angle droit : l’équerre dépasse.
\label{fig-cor-13-14}\figureGeo[corrige]{aa-14}
\end{corrige}\end{blocLivre}
\begin{blocLivre}{cor-13-15}\begin{corrige}{15}{Grand tour, petite aire}\demarche{Mesurer les deux grandeurs sur chaque figure.}
Le serpentin a un contour de $\res{20}$ unités, le carré de $\res{12}$ unités seulement. Pourtant chacun couvre $\res9$ carreaux : même aire ! L’affirmation est $\res{\text{fausse}}$ : contour et surface sont des grandeurs indépendantes.
\label{fig-cor-13-15}\figureGeo[corrige]{aa-15}
\end{corrige}\end{blocLivre}
